\section{课程安排与核心内容}
\subsection{核心内容与学习目标}
授课日期为 2026 年 9 月 17 日。\term{卷积}{convolution}是信号、图像和视频处理中的常见并行计算模式，也与科学计算中的\term{模板计算}{stencil computation}及神经网络有关。学习目标是理解卷积，并利用缓存提高 GPU 数据访问效率。一维与二维邻域加权和确定了输入访问模式，常量缓存与共享内存分别利用系数一致性和输入窗口重叠提高数据复用。

\subsection{课程安排与阅读材料}
以下安排对应上述授课日期。

\begin{table}[H]
\centering\small
\begin{tabularx}{\textwidth}{@{}p{0.23\textwidth}Y@{}}
\toprule
事项 & 安排 \\
\midrule
Lab 2 & 授课当周的星期五截止。\\
Midterm 1 时间 & 2026 年 10 月 6 日，19:00—22:00。\\
地点与形式 & 具体考场分配发布在 Canvas；纸笔考试。\\
考试范围 & Lectures 1—11，Labs 1—4。\\
时间冲突安排 & 有有效时间冲突者须在 2026 年 9 月 30 日前邮件联系教师，另行安排考试。\\
复习与阅读 & 复习资料位于 Canvas 的 Files/MT1；指定阅读为教材第 7 章。\\
\bottomrule
\end{tabularx}
\caption*{课程安排与阅读材料。}
\end{table}

\begin{keybox}[title={核心关系}]
每个输出依赖一个输入邻域，而不同输出使用相同的滤波器。前一种重叠带来输入数据的复用机会，后一种一致性适合常量存储器。分块时，输出覆盖范围确定了需要载入的输入核心区与两侧 halo；线程如何分担这些载入工作，决定了三种策略的区别。
\end{keybox}

